\section{Realism 与 Validity：真实、有用、可信是三件事}

Realism，也叫 ecological validity，问 evaluation 是否捕捉真实使用。GPQA 这类考试远离真实工作；Chatbot Arena prompt 来自真实用户但分布不可控；GDPVal 和 MedHELM 试图从专业工作场景中构造任务，但隐私和可复现性会变难。

\begin{figure}[H]
\centering
\includegraphics[width=0.86\textwidth]{gdpval.png}
\caption{GDPVal：OpenAI 从 GDP 高贡献行业采样 44 个职业任务，并由资深专业人士构造。}
\figsource{CS336 Lecture 12 官方源引用图。}
\end{figure}

\begin{knowledgebox}{读图：GDPVal 的 ecological validity}
GDPVal 更接近“模型能否替代或辅助真实工作”，但也更难标准化。真实职业任务通常没有简单唯一答案，评分 rubrics 更复杂，数据也更可能私有。这里的 tradeoff 是 realism 换来了复现难度。
\end{knowledgebox}

\begin{figure}[H]
\centering
\includegraphics[width=0.9\textwidth]{medhelm-overview.png}
\caption{MedHELM：121 个临床任务，由 29 名 clinicians 提供，结合私有和公开数据。}
\figsource{CRFM MedHELM overview，本地化后供 CS336 Lecture 12 使用。}
\end{figure}

\begin{knowledgebox}{读图：医学评价为什么要离开标准考试}
标准医学考试测知识，但临床工作还包括文档、诊断上下文、沟通、风险权衡和机构流程。MedHELM 用 clinicians 提供的任务提高 realism，但也引入隐私和数据访问限制。
\end{knowledgebox}

\begin{figure}[H]
\centering
\includegraphics[width=0.88\textwidth]{clio-table4.png}
\caption{Clio：用模型分析真实用户数据，只分享聚合模式。}
\figsource{CS336 Lecture 12 官方源引用图。}
\end{figure}

\begin{warningbox}{realism 和 privacy 的张力}
越真实的 eval 越可能含有私有数据、商业流程或用户隐私；越公开可复现的 eval 越容易变得不真实或被污染。高质量评测常需要在公开性、隐私、可重复性和现实性之间折中。
\end{warningbox}

\begin{importantbox}{课堂提示：真实 eval 往往不可完全公开}
老师在这一段的隐含判断是：公开 leaderboard 和真实产品评估之间天然有缝。公开 eval 需要可复现、可比较、可审计；真实 eval 需要贴近用户、业务和私有数据。越靠近真实使用，越可能受隐私、商业机密和数据授权限制。一个成熟团队通常会同时维护公开 benchmark、fresh eval、private eval 和线上反馈，而不是押注单一榜单。
\end{importantbox}

Validity 则问：我们如何知道 eval 真的有效？最常见威胁是 train-test overlap。传统 ML 时代 train/test split 清楚；foundation model 时代训练数据常来自互联网，评测集可能已经出现过。

\begin{figure}[H]
\centering
\includegraphics[width=0.78\textwidth]{contamination-exchangeability.png}
\caption{通过 exchangeability 线索推断 train-test overlap。}
\figsource{CS336 Lecture 12 官方源引用图。}
\end{figure}

\begin{knowledgebox}{读图：foundation model 时代的污染问题}
模型在互联网大规模训练，提供商通常不公开数据。评测集可能被训练过，导致 benchmark 高估能力。污染检测可以从模型行为推断，也可以靠报告规范、新鲜 eval 或私有 eval 降低风险。
\end{knowledgebox}

这里的 Train-test overlap 不是一个小洁癖问题，而是 foundation model 评测的基本信任危机。预训练语料覆盖互联网，benchmark 题目也在互联网传播；模型供应商又不公开完整训练集。于是我们很难证明“模型会做这个题”是泛化能力，还是见过原题或近似答案。课程列出的四条路线各有代价：从模型行为推断 overlap 依赖统计假设；报告规范依赖供应商诚实；fresh eval 会被时间戳和复制传播破坏；private eval 更真实但难以公开复现。

Dataset quality 是另一个 validity 威胁。测试用例不足、答案错误、任务太 trivial，都会让 benchmark 产生假信号。SWE-Bench Verified、Platinum benchmarks 和 Docent 都在修复“评测本身不可靠”的问题。

\begin{figure}[H]
\centering
\includegraphics[width=0.9\textwidth]{platinum-1.jpg}
\caption{Platinum benchmark 示例之一：更严格检查题目和测试质量。}
\figsource{课程源引用社交媒体图片，本地化后供 CS336 Lecture 12 使用。}
\end{figure}

\begin{figure}[H]
\centering
\includegraphics[width=0.9\textwidth]{platinum-2.jpg}
\caption{Platinum benchmark 示例之二：强调 benchmark 本身也要被审计。}
\figsource{课程源引用社交媒体图片，本地化后供 CS336 Lecture 12 使用。}
\end{figure}

\begin{warningbox}{读图：benchmark bugs 会变成模型能力幻觉}
如果测试用例不足、答案错误、任务太 trivial，agent 可能用投机策略通过，而不是具备真实能力。评价质量本身需要评价；否则模型开发会追逐坏指标。
\end{warningbox}

\begin{knowledgebox}{讲义提醒：如何审计一个 benchmark}
一份 benchmark 至少要问六个问题：样本是否代表目标任务？答案是否可靠？是否存在训练污染？题目是否已经被 frontier models 饱和？评分器是否有偏？失败样例能否指导模型改进？如果这些问题没有答案，benchmark 分数就只能作为粗略信号，而不能作为发布、采购或研究结论的唯一依据。
\end{knowledgebox}

